% chapter: wiki_vna
\chapter{VNA Measurements and Calibration}\label{ch:wiki_vna}

A vector network analyser (VNA) measures the complex S-parameters of a device under test (DUT) — both magnitude and phase — across a swept frequency range. Understanding what it measures, and how calibration moves the reference plane to your device, is essential for making accurate RF measurements.

      

\section{What a VNA Measures}

      

A 2-port VNA applies a stimulus at port 1 and measures the signals at both ports simultaneously. It computes:

\rffig{wiki_vna_1.pdf}{VNA architecture: swept source, directional couplers separate incident/reflected/transmitted waves, receivers measure amplitude and phase.}

      
        
        
        
        
        
      

{\small\begin{longtable}{>{\raggedright\arraybackslash}p{0.307\linewidth}>{\raggedright\arraybackslash}p{0.307\linewidth}>{\raggedright\arraybackslash}p{0.307\linewidth}}
\hline
\textbf{Parameter} & \textbf{Physical meaning} & \textbf{Typical display} \\
\hline
S₁₁ & Input reflection coefficient. How much is reflected from port 1. & Smith Chart, log mag (RL) \\
S₂₁ & Forward transmission. Gain or insertion loss from port 1 to port 2. & Log magnitude (dB) \\
S₁₂ & Reverse transmission. Isolation. & Log magnitude (dB) \\
S₂₂ & Output reflection coefficient. & Smith Chart, log mag \\
\hline
\end{longtable}}

      

\section{Calibration: Moving the Reference Plane}

      

Without calibration, a VNA measures its own cables, connectors, and switch imperfections. Calibration removes these systematic errors by measuring known standards and computing correction factors. The most common calibration is \textbf{SOL} (Short-Open-Load) or \textbf{SOLT} (Short-Open-Load-Through):

\rffig{wiki_vna_2.pdf}{SOL calibration standards place known Γ values on the Smith Chart rim and centre, allowing systematic error correction.}

      
\begin{itemize}

        \item \textbf{Short}: Γ = −1 (all reflected, 180° phase). Defines the reference plane position.

        \item \textbf{Open}: Γ = +1 (all reflected, 0° phase, adjusted for fringing capacitance).

        \item \textbf{Load}: Γ = 0 (50 Ω termination, no reflection). Sets the amplitude reference.

        \item \textbf{Through}: Connects port 1 to port 2 directly. Characterises transmission path.

      
\end{itemize}

      

Always calibrate at the ends of the test cables, not at the VNA ports — calibration moves the reference plane to wherever you connect the standards.

      

\section{Port Extension}

      

If you cannot calibrate at the device terminals (e.g. a connector is on a PCB and you calibrate at the cable end), port extension electrically "moves" the reference plane by adding a known electrical delay. It rotates the phase of S₁₁ and S₂₂ to compensate for the cable length. It is an approximation — it assumes the line is lossless and dispersionless — but is adequate for many applications.

      

\section{Time-Domain Gating}

      

The VNA can apply an inverse Fourier transform to its frequency-domain data to show the time-domain response (like a TDR). This reveals reflections from connectors and discontinuities separated in time (i.e. distance). Gating in the time domain — windowing to select only a specific time region — then transforming back to frequency allows you to isolate the response of a specific section of the DUT and remove connector effects. This is particularly useful for antenna measurements and filter characterisation.

      

\section{Interpreting S₁₁ on a Smith Chart}

      

A single frequency point on the Smith Chart shows the impedance at that frequency. As frequency sweeps:

      
\begin{itemize}

        \item A resistor traces a point on the real axis.

        \item A capacitor traces clockwise along the X = 0 boundary toward the short-circuit point.

        \item A series LC traces a loop: inductive (upper half), resonant (real axis), capacitive (lower half).

        \item A transmission line section traces a clockwise circle centred on the origin.

        \item An antenna shows a small loop near the edge of the chart, touching or nearly touching the real axis at resonance.

      
\end{itemize}

      

\section{Practical Measurement Tips}

      
\begin{itemize}

        \item Torque connectors to spec (SMA: 0.9 N·m / 8 in·lb). Under/over-torquing causes repeatable but incorrect results.

        \item Let the VNA warm up for 15–30 minutes before calibrating for best stability.

        \item Calibrate with the cable attached in its final position — bending it after calibration introduces errors.

        \item Use at least 201 frequency points across your band; use 1601+ for time-domain measurements to maximise time resolution.

        \item When measuring low-loss filters (\textless{} 0.5 dB), use a full 4-port SOLR or TRL calibration for best accuracy.

        \item Check calibration by measuring a known open or short — it should show a perfect circle on the Smith Chart rim.

      
\end{itemize}
